\section{Introducción}

El presente capítulo define la arquitectura del sistema inteligente de gestión
de parqueaderos, traduciendo los fundamentos teóricos desarrollados en el
capítulo anterior en decisiones concretas de diseño. Mientras el marco teórico
estableció \textit{qué} técnicas de visión artificial y aprendizaje profundo
existen y \textit{por qué} resultan adecuadas para el problema planteado, este
capítulo responde a las preguntas de \textit{cómo} se organizan dichas técnicas
en una solución integrada, \textit{qué} componentes la conforman y \textit{cómo}
interactúan entre sí para automatizar la identificación vehicular y el registro
de eventos de ingreso y salida.

El diseño se estructura de forma descendente (\textit{top-down}), partiendo de
una visión global del sistema y refinando progresivamente cada uno de sus
subsistemas. En primer lugar, se especifican los requerimientos funcionales, no
funcionales y de infraestructura que delimitan el alcance del sistema y sirven
como criterio de aceptación para las etapas posteriores de implementación y
evaluación. A continuación, se describe la arquitectura general, identificando
los componentes principales y las interfaces de comunicación entre el punto de
captura, el motor de percepción y la plataforma de gestión.

Posteriormente, el capítulo profundiza en el diseño del subsistema de visión
artificial, que constituye el núcleo del proyecto y articula el \textit{pipeline}
de detección, reconocimiento y seguimiento vehicular descrito en el marco
teórico. Luego se aborda el diseño del subsistema de gestión, compuesto por la
interfaz de programación de aplicaciones (API), el modelo de datos y los
mecanismos de persistencia de la información. Finalmente, se presenta el diseño
de la interfaz de usuario orientada al operador del recinto.

Es importante destacar que este capítulo se limita a la fase de diseño y
modelado: no incluye código fuente ni resultados experimentales, los cuales
corresponden al capítulo de implementación y pruebas. Las decisiones aquí
documentadas, tales como los umbrales de desempeño, los diagramas de
arquitectura y los modelos de datos, constituyen la referencia formal contra la
cual se verificará el sistema construido.
